\chapter{Corrected Aubrun bidefect route}

This blueprint tracks the formal proof route for ticket U-AUB-02.  The goal is
to close the corrected Aubrun count theorem used by the expectation pipeline.
The current route is not the old tight-fiber route and not the false
total-defect route.  It is the one-sided plus-defect route, converted to a
bidefect count and then absorbed by the shifted-base scalar estimate.

\section{Bidefect endpoint}

\begin{theorem}[Conditional gap-three Biane frontier with child primitive absorption]
  \label{thm:large-gap-minimal-frontier-endpoint}
  \lean{PptFactorization.AppendixB.AubrunBiDefectCountBound_of_bianeAdjacentLinkHalves_left_gap_three_start_impossible_frontiers_nestedStrictReturnIntervalOpenNonwrapOffCycleSmallerInternalOnly_and_chapuy}
  \leanok
  \uses{thm:gap-two-biane-left, thm:gap-two-biane-right, thm:left-gap-three-fixed-a1, thm:left-gap-three-S-a1-c, thm:right-gap-three-S-b-c, thm:right-gap-three-T-b1-c, thm:right-gap-three-fixed-b1, thm:right-gap-three-S-b1-d, thm:skip-branch-normalized-return, thm:skip-strict-return-frontier, thm:skip-residual-return-frontier, thm:skip-subinterval-descended-residual, thm:smaller-residual-descended-branch-cases, thm:smaller-child-primitive-absorbed}
  This is not a no-input proof of U-AUB-02.  Fix a Wick length \(m\).
  Assume the remaining large-gap left and right
  Biane adjacent-link statements outside the checked gap-three slices,
  assume the side-tied residual branch collision statements: the
  primitive-residual branches for nonwrapping first return, wrap-below last
  return, and wrap-above first return; the internal-subinterval branches for
  wrap-below and wrap-above; and the nonwrapping off-cycle open-internal
  branch.
  In the nonwrapping internal-subinterval branch, the checked
  first-return refinement removes both the endpoint-touching subcase and the
  subcase where the internal witness remains in the skipped \(x\)-cycle; the
  checked noncrossing confinement lemma below then traps the whole internal
  witness block strictly inside the return interval.
  Also assume the Chapuy recurrence
  \[
    N_{m,g+1} \le (2(m+1))^3 N_{m,g}.
  \]
  Then the Aubrun bidefect count bound holds at length \(m\).
  The phrase ``checked endpoint'' here means checked implication from these
  hypotheses.  It does not mean that the Biane, side-tied return-interval, or
  Chapuy hypotheses have been proved.  This is the current honest theorem face:
  it keeps the predecessor-skip origin data, still exposes the shorter-interval
  alternative as an actual internal subinterval of the side-tied return
  interval, splits the side disjunction into three theorem-facing origin-tag
  cases, splits each origin-tag case into primitive-residual versus
  internal-subinterval branches, exposes the nonwrapping internal-subinterval
  leaf only in its off-cycle open-internal form, sends that off-cycle witness
  to a smaller side-aware child interval, splits the child into primitive and
  internal-subinterval branches, and then absorbs the primitive child branch
  into the already theorem-facing parent primitive residual leaf.  The only
  remaining smaller child input is the genuine internal-subinterval branch.
  The loose
  primitive endpoint is kept only as a conditional diagnostic; a checked no-go
  theorem below shows that its primitive residual input is false without the
  nested side data.
  The right gap-three branch theorems listed as dependencies are checked
  internal dischargers; they are not hypotheses of the endpoint statement.
  In particular, the active endpoint no longer exposes the right special-slice
  parameters corresponding to \(S(b,c)\), \(b+1\sim_\pi c\), or
  \(\pi(b+1)=b+1\), and it no longer exposes the final branch
  \(S(b+1,d)\).  It also no longer exposes the left end-slice branch
  \(S(a+1,c)\).  The checked finite gap-three branch list is closed at this
  endpoint.
\end{theorem}

\begin{proof}
  \leanok
  The checked gap-two Biane theorems supply the adjacent links when the visible
  gap has length two.  The checked gap-three reductions remove both left
  slices \(c-a=3,b=a+1\) and \(c-a=3,b=a+2\), together with the right slice
  \(d-b=3,c=b+2\), from the broad large-gap hypotheses.  The left
  start-adjacent slice \(c-a=3,b=a+1\) is now fully discharged at this
  endpoint: the \(S(b,c)\) branch is supplied from the right frontier, the
  fixed-\(a+2\) and \(\pi(b,a+2)\) branches force that same contradiction, and
  the residual \(\neg(a\sim_\pi a+2)\wedge S(a,a+2)\) branch is contradicted
  by tree uniqueness for a copied length-two contour path.  In the left
  end-adjacent slice, the fixed \(a+1\) branch is also closed by comparing the
  copied length-four path \(S(a)\to S(c)\) with the direct path from
  \(a\sim_\pi c\).  The remaining branch \(S(a+1,c)\) forces
  \(a+1\sim_\pi b\) by comparing the copied contour step \(a+1\to b\) with the
  reversed contour step \(b\to c\), and the already checked
  \(a+1\sim_\pi b\) branch absorber closes it;
  in the right slice, the \(\pi(b,c)\) branch is closed by the immediate
  contradiction with the separating non-link \(a\not\sim_\pi b\), and the
  \(S(b,c)\) branch is closed by copying the contour step \(c\to d\) across
  \(S(b)=S(c)\) and comparing it with the direct path from
  \(b\sim_\pi d\).  The \(\pi(b+1,c)\) branch is routed through the shifted
  crossing \((a,b,b+1,d)\): it gives \(a\sim_\pi b+1\), the large-gap right
  input gives \(S(b+1,d)\).  The fixed branch \(\pi(b+1)=b+1\) is also routed
  through \(S(b+1,d)\): fixedness gives the adjacent repeat \(S(b+1,c)\), and
  denying \(S(b+1,d)\) creates a length-four incidence path from \(S(b)\) to
  \(S(d)\) that contradicts the direct length-two path from \(b\sim_\pi d\).
  Finally, the \(S(b+1,d)\) branch itself is closed: denying \(S(c,d)\) forces
  \(b+1\sim_\pi c\) by path uniqueness, the shifted left adjacent link gives
  \(S(a,b)\), and the resulting four incidences form a forbidden square.
  These links give noncrossing.  The direct skip-hit case is consumed
  internally by the checked theorem
  \(\texttt{permCyclicPredecessorSkip\_direct\_commonCycle}\).  Thus the
  theorem-facing return-interval theorem only has to handle the nested
  non-direct case.  The branch-normalized return data is supplied internally,
  the full-endpoint common-cycle branch is discharged internally, and the raw
  internal right-repeat branch is first replaced by a subinterval-preserving
  descended pair of alternatives: an adjacent right repeat or a strictly
  shorter longer same-\(\pi\) subinterval \([a+r,a+s]\) inside the current
  return interval.  The theorem-facing return-collision input is therefore the
  side-tied subinterval-descended residual frontier, not the loose primitive
  five-branch frontier.  The separate primitive descent theorem is only a
  conditional diagnostic unless the predecessor-skip side data is carried
  through descent.  In the nonwrapping off-cycle child interval, the checked
  reductions either find the common pair immediately or preserve the side data
  on a smaller child residual; the primitive child part of that residual is
  absorbed back into the parent primitive residual leaf.  Thus the remaining
  smaller child theorem-facing branch is only the genuine
  internal-subinterval branch.
  Together with left-tight transversality this rules out
  non-fixed predecessor skips.  The existing genus-zero and Chapuy route then
  gives the bidefect count bound.  This is a conditional reduction theorem:
  the \(\leanok\)
  marker certifies the reduction from the named hypotheses, not the open
  large-gap Biane, side-tied subinterval return-interval, or Chapuy inputs
  themselves.
\end{proof}

\begin{theorem}[Off-cycle internal blocks are strictly nested]
  \label{thm:off-cycle-block-confinement}
  \lean{PptFactorization.AppendixB.leftTightStrictReturnIntervalOpenInternalOffCycleSubintervalResidual_block_strictly_inside}
  \leanok
  In the nonwrapping first-return branch, suppose \(a\) and \(c\) lie in the
  same left cycle, the transported cycle partition is noncrossing, and the
  internal subinterval witness starts at a point \(a+r\) which is not in the
  skipped \(x\)-cycle.  Then every point of the left cycle containing \(a+r\)
  lies strictly between \(a\) and \(c\).
\end{theorem}

\begin{proof}
  \leanok
  The points \(a\) and \(c\) lie in the skipped \(x\)-cycle, while \(a+r\) is
  off that cycle.  Thus the block of \(a+r\) is distinct from the block of
  \(a,c\), and it contains a point strictly between \(a\) and \(c\).  The
  noncrossing interval-nesting lemma confines the whole block to the closed
  interval \([a,c]\).  If one of its points were equal to \(a\) or \(c\), the
  block would meet the skipped \(x\)-cycle, contradicting the off-cycle
  condition.  Therefore the confinement is strict.
\end{proof}

\begin{theorem}[Off-cycle witnesses give smaller contour intervals]
  \label{thm:off-cycle-smaller-contour-interval}
  \lean{PptFactorization.AppendixB.leftTightStrictReturnIntervalOpenInternalOffCycleSubintervalResidual_strictSubinterval_contourRepeat}
  \leanok
  In the same nonwrapping off-cycle branch, the internal witness interval
  \([a+r,a+s]\) is itself a strictly smaller same-\(\pi\) return interval
  inside \((a,c)\).  It carries the standard
  \texttt{leftTightContourIntervalRepeat} obstruction, and every point of its
  \(\pi\)-block remains off the skipped \(x\)-cycle.
\end{theorem}

\begin{proof}
  \leanok
  The endpoints \(a+r\) and \(a+s\) come directly from the off-cycle internal
  residual witness.  The inequalities \(0<r<s<c-a\) and \(1<s-r\) make this a
  strictly smaller nontrivial interval.  The previous confinement theorem and
  the first-return condition show that the whole witness block lies in
  \((a,c)\) and is disjoint from the skipped \(x\)-cycle.  Finally the generic
  left-tight contour-repeat theorem applies to the longer same-\(\pi\)
  interval \([a+r,a+s]\).
\end{proof}

\begin{theorem}[Off-cycle witnesses reduce to smaller side-aware residuals]
  \label{thm:off-cycle-smaller-side-aware-residual}
  \lean{PptFactorization.AppendixB.leftTightStrictReturnIntervalOpenInternalOffCycleSubintervalResidual_common_or_smallerResidual}
  \leanok
  In the nonwrapping off-cycle branch, the smaller contour-repeat interval
  either already gives a common \(\pi\)/\(S\pi\) pair, or leaves a strict
  residual on a smaller interval.  In the latter case, no point of the skipped
  \(x\)-cycle lies inside the child interval, and the child \(\pi\)-block
  remains inside the parent interval and off the skipped cycle.
\end{theorem}

\begin{proof}
  \leanok
  Apply the smaller-contour theorem, then split the resulting contour repeat by
  the strict-frontier decomposition.  The common branch is immediate.  In the
  residual branch, the first-return condition restricts from the parent
  interval to the child interval, and the block confinement data is carried
  unchanged.
\end{proof}

\begin{theorem}[The smaller residual splits into descended child branches]
  \label{thm:smaller-residual-descended-branch-cases}
  \lean{PptFactorization.AppendixB.leftTight_nonwrap_smallerResidual_common_of_descendedBranchCases}
  \leanok
  The side-aware residual child from the nonwrapping off-cycle branch can be
  decomposed by the subinterval-preserving descent.  Thus the theorem-facing
  child input is not a raw residual predicate: it is either the side-aware
  primitive residual branch or the side-aware internal-subinterval branch,
  with the no-\(x\)-cycle and block-off-\(x\) data still available.
\end{theorem}

\begin{proof}
  \leanok
  Apply the checked descent from a raw strict residual to the
  subinterval-preserving descended residual.  The definitional branch split of
  that descended residual gives the primitive child case or the internal
  subinterval child case, and each case is passed to the corresponding
  side-aware collision input.
\end{proof}

\begin{theorem}[The smaller primitive child splits into five local branches]
  \label{thm:smaller-primitive-child-branch-cases}
  \lean{PptFactorization.AppendixB.leftTight_nonwrap_smallerPrimitive_common_of_branchCases}
  \leanok
  The primitive child residual in the nonwrapping off-cycle branch admits a
  checked five-way local split:
  internal \(S\)-repeat, start-attached \(S\)-repeat, end-attached
  \(S\)-repeat, fixed predecessor, and adjacent \(T\)-repeat.  Each case still
  carries the outer nonwrapping first-return data, the no-\(x\)-cycle condition
  inside the child interval, and the block-off-\(x\) invariant.  The current
  endpoint no longer exposes these five cases separately; the following
  arithmetic bridge absorbs them into the parent primitive leaf.
\end{theorem}

\begin{proof}
  \leanok
  Unfold the primitive residual frontier into its definitional five-way branch
  split.  Each branch is passed to the corresponding side-aware child
  collision input with the inherited first-return, no-\(x\)-cycle, and
  block-off-\(x\) hypotheses unchanged.
\end{proof}

\begin{theorem}[Strict child primitive residuals are parent primitive residuals]
  \label{thm:smaller-child-primitive-absorbed}
  \lean{PptFactorization.AppendixB.leftTightStrictReturnIntervalPrimitiveResidual_of_strictChildPrimitive}
  \leanok
  Let \([a',c']\) be a strict child interval inside a parent interval
  \([a,c]\).  If the child interval carries any primitive residual branch, then
  the parent interval itself carries a primitive residual branch.
\end{theorem}

\begin{proof}
  \leanok
  The proof is arithmetic.  Internal, start-attached, and end-attached child
  \(S\)-repeats become internal parent \(S\)-repeats after translating the
  offsets from \(a'\) to \(a\).  A fixed predecessor at \(c'-1\) gives the
  adjacent left repeat \((c'-1,c')\), again an internal parent \(S\)-repeat.
  A child adjacent \(T\)-repeat remains an adjacent \(T\)-repeat after the same
  offset translation.  Thus no new child primitive collision theorem is
  needed beyond the parent primitive residual leaf.
\end{proof}

\section{Checked local Biane inputs}

\begin{theorem}[Left gap-two Biane link]
  \label{thm:gap-two-biane-left}
  \lean{PptFactorization.AppendixB.leftTight_bianeAdjacentLink_left_of_gap_two}
  \leanok
  For a crossing \(a<b<c<d\) with \(c-a=2\), the left adjacent
  \((\gamma\pi)\)-cycle link \(a\sim_{\gamma\pi} b\) follows from
  left-tightness, the crossing links \(a\sim_\pi c\), \(b\sim_\pi d\), and the
  separating non-link \(a\not\sim_\pi b\).
\end{theorem}

\begin{theorem}[Right gap-two Biane link]
  \label{thm:gap-two-biane-right}
  \lean{PptFactorization.AppendixB.leftTight_bianeAdjacentLink_right_of_gap_two}
  \leanok
  For a crossing \(a<b<c<d\) with \(d-b=2\), the right adjacent
  \((\gamma\pi)\)-cycle link \(c\sim_{\gamma\pi} d\) follows from the same
  left-tight crossing data.
\end{theorem}

\begin{theorem}[Left gap-three start-adjacent reduction]
  \label{thm:left-gap-three-reduction}
  \lean{PptFactorization.AppendixB.leftTight_bianeAdjacentLink_left_of_gap_three_start_adjacent_reduced_fixed_a2_no_pi_a_a2}
  \leanok
  In the finite slice \(c-a=3\), \(b=a+1\), the contour-repeat split first
  reduces the left Biane link to four explicit start-slice branches:
  the residual branch \(\neg(a\sim_\pi a+2)\) together with \(S(a,a+2)\),
  \(S(b,c)\), the fixed-point branch \(\pi(a+2)=a+2\), and \(\pi(b,a+2)\).
  Later checked suppliers discharge all four at the endpoint, so this reducer
  is now historical plumbing for the fully closed left start-adjacent slice.
\end{theorem}

\begin{theorem}[Left start target impossibility]
  \label{thm:left-start-target-impossible}
  \lean{PptFactorization.AppendixB.leftTight_bianeAdjacentLink_left_of_gap_three_start_adjacent_branch_to_target_iff_false}
  \leanok
  In the same finite slice \(c-a=3\), \(b=a+1\), the target \(S(a,b)\) is
  incompatible with the nontrivial \(\pi\)-interval \(a\sim_\pi c\).  Hence
  each remaining branch-to-target obligation in this start slice is equivalent
  to proving that the branch itself is impossible.
\end{theorem}

\begin{theorem}[Residual left-start branch]
  \label{thm:left-start-residual}
  \lean{PptFactorization.AppendixB.leftTight_bianeAdjacentLink_left_of_gap_three_start_S_a_a2_no_pi_false}
  \leanok
  In the slice \(c-a=3\), \(b=a+1\), the branch
  \(\neg(a\sim_\pi a+2)\wedge S(a,a+2)\) is impossible.  The repeat
  \(S(a,a+2)\) copies the contour step from \(a+2\) to \(c\) into a
  length-two path from \(S(a)\) to \(S(c)\); comparison with the direct
  length-two path from \(a\sim_\pi c\) forces \(a\sim_\pi a+2\).
\end{theorem}

\begin{theorem}[Left gap-three end-adjacent reduction]
  \label{thm:left-gap-three-end-reduction}
  \lean{PptFactorization.AppendixB.leftTight_bianeAdjacentLink_left_of_gap_three_end_adjacent_reduced_fixed_a1_no_pi_branches}
  \leanok
  In the finite slice \(c-a=3\), \(b=a+2\), the contour-repeat split reduces
  the left Biane link to two explicit branch implications.  The adjacent
  repeat \(S(a+1,b)\) has been normalized, using left-tightness, to the
  fixed-point branch \(\pi(a+1)=a+1\).  The other remaining branch is
  \(S(a+1,c)\).  The next two theorems close both branches, so this finite
  slice is no longer theorem-facing on the sharp endpoint.
\end{theorem}

\begin{theorem}[Left gap-three fixed \(a+1\) branch]
  \label{thm:left-gap-three-fixed-a1}
  \lean{PptFactorization.AppendixB.leftTight_bianeAdjacentLink_left_of_gap_three_end_adjacent_fixed_a1}
  \leanok
  In the same slice \(c-a=3\), \(b=a+2\), fixedness of \(a+1\) gives the
  adjacent repeat \(S(a+1,b)\).  If the target \(S(a,b)\) failed, the contour
  step \(a\to a+1\), this repeat, and the contour step \(b\to c\) would form a
  simple length-four incidence path from \(S(a)\) to \(S(c)\).  Tree
  uniqueness compares it with the direct length-two path from \(a\sim_\pi c\),
  contradiction.
\end{theorem}

\begin{theorem}[Left gap-three \(S(a+1,c)\) branch]
  \label{thm:left-gap-three-S-a1-c}
  \lean{PptFactorization.AppendixB.leftTight_bianeAdjacentLink_left_of_gap_three_end_adjacent_T_a1_b_of_S_a1_c}
  \leanok
  In the same slice \(c-a=3\), \(b=a+2\), the branch \(S(a+1,c)\) forces
  \(a+1\sim_\pi b\).  The copied contour step \(a+1\to b\) is a simple path
  from \(S(c)\) to \(S(b)\), and the reversed contour step \(b\to c\) is
  another simple path with the same endpoints.  Tree uniqueness identifies
  their middle right vertices.  The already checked
  \(a+1\sim_\pi b\) branch theorem then closes the target.
\end{theorem}

\begin{theorem}[The \(\pi(a,a+1)\) end-adjacent branch]
  \label{thm:left-gap-three-pi-a-a1}
  \lean{PptFactorization.AppendixB.leftTight_bianeAdjacentLink_left_of_gap_three_end_adjacent_of_pi_a_a1}
  \leanok
  In the same slice, the branch \(\pi(a,a+1)\) is absorbed by the shifted
  gap-two theorem: \(\pi(a,a+1)\) and \(\pi(a,c)\) imply
  \(\pi(a+1,c)\), so the crossing \((a+1,b,c,d)\) gives \(S(a+1,b)\), which
  is one of the remaining explicit branches.
\end{theorem}

\begin{theorem}[The \(\pi(a+1,b)\) end-adjacent branch]
  \label{thm:left-gap-three-pi-a1-b}
  \lean{PptFactorization.AppendixB.leftTight_bianeAdjacentLink_left_of_gap_three_end_adjacent_of_pi_a1_b}
  \leanok
  In the same slice, the branch \(\pi(a+1,b)\) is absorbed by the already
  exposed start-adjacent branch on the shifted crossing \((a,a+1,c,d)\).
  That branch gives the adjacent repeat \(S(a,a+1)\), which contradicts the
  nontrivial same-\(\pi\) interval \(a\sim_\pi c\).
\end{theorem}

\begin{theorem}[Right gap-three end-adjacent reduction]
  \label{thm:right-gap-three-reduction}
  \lean{PptFactorization.AppendixB.leftTight_bianeAdjacentLink_right_of_gap_three_end_adjacent_reduced_fixed_b1_no_pi_b_b1_no_pi_b_c}
  \leanok
  In the finite slice \(d-b=3\), \(c=b+2\), the contour-repeat split and the
  shifted gap-two branch remove the \(\pi(b,b+1)\) case, while transitivity
  with \(a\sim_\pi c\) removes the \(\pi(b,c)\) case.  The remaining branches
  at this reducer level are \(S(b,c)\), \(\pi(b+1)=b+1\), \(S(b+1,d)\), and
  \(\pi(b+1,c)\).  The adjacent repeat \(S(b+1,c)\) has been normalized,
  using left-tightness, to the fixed-point branch.  The separate theorem below
  closes the \(S(b,c)\) branch at the current endpoint, and the next theorem
  routes the \(\pi(b+1,c)\) branch through \(S(b+1,d)\).  The fixed-branch
  theorem after that routes \(\pi(b+1)=b+1\) through the same \(S(b+1,d)\)
  branch, and the final theorem below closes \(S(b+1,d)\) itself.  Thus the
  active endpoint has no theorem-facing right branch in this checked slice.
\end{theorem}

\begin{theorem}[Right gap-three \(S(b,c)\) branch]
  \label{thm:right-gap-three-S-b-c}
  \lean{PptFactorization.AppendixB.leftTight_bianeAdjacentLink_right_of_gap_three_end_adjacent_of_S_b_c}
  \leanok
  In the finite slice \(d-b=3\), \(c=b+2\), the branch \(S(b,c)\) forces the
  target \(S(c,d)\).  If \(S(c,d)\) failed, the contour step from \(c\) to
  \(d\), copied across \(S(b)=S(c)\), would be a simple path from \(S(b)\) to
  \(S(d)\).  Tree uniqueness compares it with the direct path from
  \(b\sim_\pi d\), forcing \(b\sim_\pi c\), which contradicts the crossing
  link \(a\sim_\pi c\) and the separating non-link \(a\not\sim_\pi b\).
\end{theorem}

\begin{theorem}[Right gap-three \(\pi(b+1,c)\) branch]
  \label{thm:right-gap-three-T-b1-c}
  \lean{PptFactorization.AppendixB.leftTight_bianeAdjacentLink_right_of_gap_three_end_adjacent_T_b1_c_to_S_b1_d_of_right_rest}
  \leanok
  In the finite slice \(d-b=3\), \(c=b+2\), the branch
  \(b+1\sim_\pi c\) gives \(a\sim_\pi b+1\) by transitivity with the crossing
  link \(a\sim_\pi c\).  Applied to the shifted crossing \((a,b,b+1,d)\), the
  large-gap right input supplies \(S(b+1,d)\), because this shifted crossing
  is not the special \(c=b+2\) slice.  The final right-branch theorem below
  consumes \(S(b+1,d)\), so \(b+1\sim_\pi c\) is no longer theorem-facing.
\end{theorem}

\begin{theorem}[Right gap-three fixed \(b+1\) branch]
  \label{thm:right-gap-three-fixed-b1}
  \lean{PptFactorization.AppendixB.leftTight_bianeAdjacentLink_right_of_gap_three_end_adjacent_fixed_b1_to_S_b1_d}
  \leanok
  In the finite slice \(d-b=3\), \(c=b+2\), the fixed branch
  \(\pi(b+1)=b+1\) gives the adjacent repeat \(S(b+1,c)\).  If \(S(b+1,d)\)
  failed, the contour step \(b\to b+1\), followed by the copied contour step
  \(c\to d\), would form a simple length-four incidence path from \(S(b)\) to
  \(S(d)\).  Tree uniqueness compares it with the direct length-two path
  coming from \(b\sim_\pi d\), contradiction.  Hence the fixed branch is no
  longer theorem-facing at the active endpoint; it is fed to the final
  \(S(b+1,d)\) branch theorem.
\end{theorem}

\begin{theorem}[Right gap-three \(S(b+1,d)\) branch]
  \label{thm:right-gap-three-S-b1-d}
  \lean{PptFactorization.AppendixB.leftTight_bianeAdjacentLink_right_of_gap_three_end_adjacent_of_S_b1_d_of_left_shift}
  \leanok
  In the finite slice \(d-b=3\), \(c=b+2\), the branch \(S(b+1,d)\) forces the
  target \(S(c,d)\), provided the shifted left adjacent link on
  \((a,b,b+1,d)\) is available.  If \(S(c,d)\) failed, uniqueness of the two
  simple incidence paths from \(S(c)\) to \(S(d)\) would force
  \(b+1\sim_\pi c\).  Then \(a\sim_\pi c\) gives \(a\sim_\pi b+1\), the
  shifted left link gives \(S(a,b)\), and the links
  \(a\sim_\pi b+1\), \(b\sim_\pi d\), \(S(a,b)\), and \(S(b+1,d)\) form a
  forbidden incidence square.
\end{theorem}

\section{Skip frontier}

\begin{theorem}[Minimal frontier extracted from a non-fixed skip]
  \label{thm:skip-minimal-frontier}
  \lean{PptFactorization.AppendixB.leftTight_nonfixed_cyclicPredecessorSkip_exists_minimal_frontier}
  \leanok
  Any non-fixed predecessor skip in a left-tight permutation produces a
  minimal long same-\(\pi\) interval.  Its contour frontier has five possible
  alternatives: an adjacent right repeat, an internal left repeat, a proper
  start-touching left repeat, a proper end-touching left repeat, or a fixed
  predecessor at the right endpoint.
\end{theorem}

\begin{theorem}[Tied contour interval extracted from a non-fixed skip]
  \label{thm:skip-tied-contour}
  \lean{PptFactorization.AppendixB.leftTight_nonfixed_cyclicPredecessorSkip_exists_tied_contourIntervalRepeat}
  \leanok
  Any non-fixed predecessor skip in a left-tight permutation produces a long
  repeated contour interval whose endpoints remember the skip witness.  In the
  nonwrapping case the interval is exactly
  \([(\pi x).\mathrm{val},x.\mathrm{val}]\); in the wrapping-below case it
  ends at \((\pi x).\mathrm{val}\) and starts below \(x\); in the
  wrapping-above case it starts at \(x.\mathrm{val}\) and ends above
  \((\pi x).\mathrm{val}\).
\end{theorem}

\begin{theorem}[First return in the nonwrapping nested skip branch]
  \label{thm:skip-nonwrap-first-return}
  \lean{PptFactorization.AppendixB.leftTight_nonwrap_nested_cyclicPredecessorSkip_exists_first_contourIntervalRepeat}
  \leanok
  In the nonwrapping branch, if \(\gamma(\pi x)\) is not in the
  \(\pi\)-cycle of \(x\), then the first same-\(\pi\) return after
  \((\pi x).\mathrm{val}\) is separated by a genuine gap.  The resulting
  interval starts at \((\pi x).\mathrm{val}\), has no earlier same-cycle
  return to \(x\), and carries the repeated-contour obstruction.
\end{theorem}

\begin{theorem}[First return in the wrap-above nested skip branch]
  \label{thm:skip-wrap-above-first-return}
  \lean{PptFactorization.AppendixB.leftTight_wrapAbove_nested_cyclicPredecessorSkip_exists_first_contourIntervalRepeat}
  \leanok
  In the wrap-above branch, the first same-\(\pi\) return above
  \((\pi x).\mathrm{val}\) is also separated from \(\pi x\) by a genuine gap
  unless the direct-hit case occurs.  The corresponding long interval starts
  at \(x.\mathrm{val}\) and carries the repeated-contour obstruction.
\end{theorem}

\begin{theorem}[Branch-normalized nested skip return]
  \label{thm:skip-branch-normalized-return}
  \lean{PptFactorization.AppendixB.leftTight_nonfixed_nested_cyclicPredecessorSkip_exists_returnContourIntervalRepeat}
  \leanok
  In a genuinely nested non-fixed predecessor skip, the canonical return
  interval is now available in every branch: first return in the nonwrapping
  branch, last return in the wrap-below branch, and first return in the
  wrap-above branch.  Each interval has a genuine gap and carries the
  repeated-contour obstruction supplied by left-tightness.
\end{theorem}

\begin{theorem}[Strict return-interval frontier]
  \label{thm:skip-strict-return-frontier}
  \lean{PptFactorization.AppendixB.leftTightContourIntervalRepeat_strict_frontier}
  \leanok
  A branch-normalized repeated-contour interval can be sharpened to a strict
  return frontier.  The full endpoint repeat is already a common-cycle
  collision.  The start-adjacent left repeat is impossible by left-tightness,
  and the end-adjacent left repeat is exactly the fixed-predecessor branch.
  The theorem-facing strict alternatives are therefore internal left repeat,
  proper start repeat, proper end repeat, fixed predecessor, or internal
  right repeat.
\end{theorem}

\begin{theorem}[Residual return-interval frontier]
  \label{thm:skip-residual-return-frontier}
  \lean{PptFactorization.AppendixB.leftTightStrictReturnIntervalFrontier_common_or_residual}
  \leanok
  The strict return frontier splits into an already-solved common-cycle branch
  and a residual frontier.  The active endpoint consumes only the residual
  frontier; the common-cycle branch is discharged internally.
\end{theorem}

\begin{theorem}[Descended residual return frontier]
  \label{thm:skip-descended-residual}
  \lean{PptFactorization.AppendixB.leftTightStrictReturnIntervalResidual_descend_internal_right}
  \leanok
  The residual internal right-repeat branch is no longer left raw.  A checked
  descent replaces it by either an adjacent same-\(\pi\) repeat or a strictly
  shorter longer same-\(\pi\) interval.  This is a narrowing of the nested
  collision input, not a proof of that collision input.
\end{theorem}

\begin{theorem}[Subinterval-descended residual return frontier]
  \label{thm:skip-subinterval-descended-residual}
  \lean{PptFactorization.AppendixB.leftTightStrictReturnIntervalResidual_descend_internal_right_subinterval}
  \leanok
  The residual internal right-repeat branch is narrowed without losing the
  side-tied return-interval geometry.  Adjacent internal right repeats become
  the adjacent primitive branch; non-adjacent internal right repeats are kept
  as an actual shorter subinterval \([a+r,a+s]\) of the current return
  interval.  Thus the theorem-facing recursive branch is no longer an
  arbitrary shorter same-\(\pi\) interval.
\end{theorem}

\begin{theorem}[Loose primitive residual no-go]
  \label{thm:loose-primitive-residual-no-go}
  \lean{PptFactorization.AppendixB.not_forall_loose_primitiveResidual_commonCycle}
  \leanok
  The primitive residual collision statement is false if stated without the
  nested predecessor-skip origin data.  Concretely, at \(m=2\), the inverse
  long cycle is left-tight and noncrossing, the interval \((0,2)\) satisfies
  the primitive residual predicate through the adjacent-\(\pi\) branch, but
  \(\gamma\pi=1\), so there is no nontrivial common \(\pi\)- and
  \(\gamma\pi\)-cycle pair.
\end{theorem}

\begin{theorem}[Primitive residual descent]
  \label{thm:skip-primitive-residual}
  \lean{PptFactorization.AppendixB.leftTightStrictReturnIntervalDescendedResidual_common_of_primitive}
  \leanok
  The strictly shorter same-\(\pi\) interval alternative in the descended
  residual frontier is closed by well-founded descent on interval length.  It
  is enough for the theorem-facing collision input to solve the five primitive
  residual branches: internal left repeat, proper start repeat, proper end
  repeat, fixed predecessor, and adjacent right repeat.  This theorem is not
  the current endpoint bridge by itself, because its primitive input does not
  carry the predecessor-skip side data.
\end{theorem}

\section{Open leaves}

\begin{theorem}[Large-gap left Biane link]
  \label{thm:open-large-gap-left}
  For every genuine crossing with \(2<c-a\) and \(c-a\ne3\), prove the left
  adjacent link \(a\sim_{\gamma\pi} b\).  This is still open.
\end{theorem}

\begin{theorem}[Large-gap right Biane link]
  \label{thm:open-large-gap-right}
  For every genuine crossing with \(2<d-b\) outside the checked slice
  \(d-b=3,\ c=b+2\), prove the right adjacent link
  \(c\sim_{\gamma\pi} d\).  This is still open.
\end{theorem}

\begin{theorem}[Nested return-interval branch collision]
  \label{thm:open-minimal-frontier-collision}
  Prove the side-tied residual branch collision statements consumed by the
  current endpoint, or replace them by a side-descending primitive theorem.
  These are the primitive-residual branches for nonwrapping first return,
  wrap-below last return, and wrap-above first return; the internal-subinterval
  branches for wrap-below and wrap-above; and the nonwrapping off-cycle
  open-internal branch.  In the nonwrapping internal-subinterval case, the
  endpoint-touching start \(r=0\) and the \(x\)-cycle internal-witness subcase
  are already discharged by the first-return condition, so the remaining leaf
  is the off-cycle open-internal case.  The theorem must carry
  the origin tag or carry equivalent context through interval descent.  The raw
  internal right-repeat branch has been replaced by adjacent-or-subinterval
  descent, and the recursive branch now remembers the internal offsets
  \([a+r,a+s]\) inside the side-tied return interval.  What remains is not the
  loose primitive frontier by itself:
  Lean gives a counterexample to that over-broad statement.  The remaining
  theorem must use the nested side/context data and a genuine
  noncrossing/incidence-tree argument.
\end{theorem}

\begin{theorem}[Chapuy recurrence]
  \label{thm:open-chapuy}
  Prove the slicing recurrence
  \[
    N_{m,g+1} \le (2(m+1))^3 N_{m,g}.
  \]
  This is still open.
\end{theorem}
